\section{Introduction}
\label{sec:intro}
Comparing RL post-training systems requires separating the algorithm name from
the implemented objective, sampling protocol, checkpoint, and evaluator.
Implementation sensitivity is already established in policy-gradient
research~\citep{engstrom2020implementation,henderson2018deep}; heterogeneous
observations alone do not identify which implementation choice caused a gap.

\ours{} is a bounded evidence audit. Its contributions are:
\begin{enumerate}
\item An objective-scoped account of why homogeneous rewards eliminate the
centered reward-contrast term in the specified noncanonical runner.
\item An inventory of the two reviewed synthetic tool-use cells and explicit
quarantines for conflicting or insufficiently identified comparisons.
\item A reporting and prospective validation contract distinguishing descriptive
reward summaries from prediction, stopping decisions, and intervention benefit.
\end{enumerate}
These contributions do not include a new diagnostic definition, a controlled
cross-library benchmark, a validated stopping rule, or an algorithm ranking.

Zero-Variance Fraction (ZVF) counts homogeneous reward groups. Its historical
complement, Gradient Utilization (GU), is retained as terminology, but measures
only the fraction of groups with nonconstant rewards. Such groups have centered
reward contrast; neither gradient magnitude nor usefulness follows. Our two
all-wrong tool-use observations have no held-out evaluation. They motivate a
reporting discipline, not a claim that the diagnostic predicts failure beyond
reward-only baselines.
